% Solutions/hints for ch19 exercises (assembled into Appendix C)

\begin{solution}{exr:ch19-which}
(a) Neither: $f$ is the constant-velocity matrix and $h$ selects the position; use the Kalman filter. (b) $h$ is nonlinear (square root and arctangent), $f$ linear; an EKF or UKF, the UKF when the bearing noise is large or the range short. (c) $f$ is nonlinear through $v\cos\psi$, $v\sin\psi$ and the division by $\omega$, $h$ linear; the EKF is adequate for mild maneuvers. (d) Both nonlinear; EKF or UKF at long range, UKF at short range, a particle filter if there are occlusions or a map. (e) With $\omega$ known and fixed, the Cartesian-velocity form $(x, \dot{x}, y, \dot{y})$ evolves by a constant matrix whose entries are $\sin(\omega\dt)/\omega$, $(1 - \cos\omega\dt)/\omega$, $\cos\omega\dt$ and $\sin\omega\dt$, and $h$ selects the position: a linear-Gaussian model, so the plain Kalman filter is exact. The nonlinearity of the coordinated-turn model comes entirely from $\omega$ being unknown and multiplied with the velocity.
\end{solution}

\begin{solution}{exr:ch19-ct-jacobian}
Write $\psi' = \psi + \omega\dt$. The first component is $x' = x + \frac{v}{\omega}(\sin\psi' - \sin\psi)$; its derivatives with respect to $x$, $y$, $v$ are $1$, $0$ and $(\sin\psi' - \sin\psi)/\omega$; with respect to $\psi$, since $\partial\psi'/\partial\psi = 1$, $\frac{v}{\omega}(\cos\psi' - \cos\psi)$; with respect to $\omega$, by the product rule on $v\omega^{-1}$ and $\partial\psi'/\partial\omega = \dt$, $-\frac{v}{\omega^2}(\sin\psi' - \sin\psi) + \frac{v\dt}{\omega}\cos\psi'$. The second component $y' = y - \frac{v}{\omega}(\cos\psi' - \cos\psi)$ gives, in the same way, $-(\cos\psi' - \cos\psi)/\omega$, $\frac{v}{\omega}(\sin\psi' - \sin\psi)$ and $\frac{v}{\omega^2}(\cos\psi' - \cos\psi) + \frac{v\dt}{\omega}\sin\psi'$. For the limits use $\sin\psi' = \sin\psi + \omega\dt\cos\psi - \tfrac12\omega^2\dt^2\sin\psi - \tfrac16\omega^3\dt^3\cos\psi + \dots$ and $\cos\psi' = \cos\psi - \omega\dt\sin\psi - \tfrac12\omega^2\dt^2\cos\psi + \tfrac16\omega^3\dt^3\sin\psi + \dots$: then $(\sin\psi' - \sin\psi)/\omega \to \dt\cos\psi$, $\frac{v}{\omega}(\cos\psi' - \cos\psi) \to -v\dt\sin\psi$, and in $F_{15}$ the terms of order $1/\omega$ cancel, leaving $-\tfrac12 v\dt^2\sin\psi + \bigO{\omega}$; the second row is analogous with $\dt\sin\psi$, $v\dt\cos\psi$ and $\tfrac12 v\dt^2\cos\psi$. The chapter code evaluates exactly these series for $\abs{\omega} < 10^{-4}$, and a central finite difference with step $10^{-5}$ agrees with \code{CoordinatedTurnModel.jacobian} to better than $10^{-5}$ at both requested turn rates.
\end{solution}

\begin{solution}{exr:ch19-weights}
With $n = 5$ and $\lambda = \alpha^2(n + \kappa) - n$: for $(1, 2, 0)$, $\lambda = 0$, the points are $\sqrt{5} \approx 2.236$ standard deviations from the mean, $W_0^{(m)} = 0$, $W_0^{(c)} = 2$, $W_i^{(m)} = 0.1$. For $(1, 2, -2)$, $\lambda = -2$, $n + \lambda = 3$, the points are at $\sqrt{3} \approx 1.732$ standard deviations, $W_0^{(m)} = -2/3$, $W_0^{(c)} = 4/3$, $W_i^{(m)} = 1/6$. For $(10^{-3}, 2, 0)$, $\lambda = 5 \cdot 10^{-6} - 5$, $n + \lambda = 5 \cdot 10^{-6}$, the points are $\sqrt{5} \cdot 10^{-3}$ standard deviations from the mean, $W_0^{(m)} = \lambda/(n + \lambda) = 1 - n/(n + \lambda) = 1 - 10^6$, $W_0^{(c)} = W_0^{(m)} + 3 - 10^{-6} \approx 3 - 10^6$, and $W_i^{(m)} = 1/(2(n + \lambda)) = 10^5$. The second and third choices have a negative $W_0^{(m)}$. In every case $W_0^{(m)} + 10\,W_i^{(m)} = 1$: $0 + 1$, $-2/3 + 10/6$ and $(1 - 10^6) + 10^6$. The danger of the third choice is cancellation: the mean is $(1 - 10^6)$ times the center image plus $10^5$ times the sum of ten images that differ from the center by about $10^{-3}$ standard deviations, that is, the difference of two numbers of size $10^6$ that must be accurate to a relative $10^{-6}$; six of the sixteen decimal digits of double precision are lost, and any rounding inside $g$ is amplified a millionfold.
\end{solution}

\begin{solution}{exr:ch19-ess}
(a) $N_{\mathrm{eff}} = 1/(4 \cdot 0.0625) = 4$; $1/(0.49 + 3 \cdot 0.01) = 1/0.52 \approx 1.92$; $1/1 = 1$. (b) Cauchy--Schwarz with the vectors $(\omega^{(1)}, \dots, \omega^{(N)})$ and $(1, \dots, 1)$ gives $1 = (\sum_i \omega^{(i)})^2 \le N \sum_i (\omega^{(i)})^2$, so $N_{\mathrm{eff}} \le N$, with equality exactly when the two vectors are proportional, that is, all weights equal; and $\sum_i (\omega^{(i)})^2 \le \max_i \omega^{(i)} \le 1$ with equality exactly when one weight is $1$. (c) Any set of $N$ particles with equal weights, wherever they are: take all particles $100\,\mathrm{m}$ from the truth before any measurement has been processed. $N_{\mathrm{eff}}$ measures how evenly the weight is spread over the particles, not whether the particles are anywhere near the truth; a filter that has lost the track and resampled has $N_{\mathrm{eff}} = N$. (d) The cumulative sums are $(0.7, 0.8, 0.9, 1.0)$ and the grid points $0.1, 0.35, 0.6, 0.85$. The first three fall into $(0, 0.7]$ and select particle~1; $0.85$ falls into $(0.8, 0.9]$ and selects particle~3. The counts $(3, 0, 1, 0)$ lie in $\{\lfloor 2.8 \rfloor, \lceil 2.8 \rceil\} = \{2, 3\}$, $\{0, 1\}$, $\{0, 1\}$, $\{0, 1\}$ as \cref{thm:ch19-systematic} requires, and average over $u$ to $(2.8, 0.4, 0.4, 0.4)$.
\end{solution}

\begin{solution}{exr:ch19-wrapping}
The update is linear in the innovation, so the naive estimate differs from the correct one by $\mat{K}_{\cdot 2}$ times the difference of the two innovations, and that difference is exactly $6.2117 - (-0.0715) = 2\pi$. The shifts are $3.21 \cdot 2\pi = +20.2\,\mathrm{m}$ in $x$, $-47.05 \cdot 2\pi = -295.6\,\mathrm{m}$ in $y$ and $2.56 \cdot 2\pi = +16.1\,\mathrm{rad}$ in $\psi$ (that is $+3.5\,\mathrm{rad}$ after wrapping): starting from the updated estimate $(-110.94, -2.13)$ of \cref{tab:ch19-ekf-trace}, the position lands near $(-90.8, -297.8)$, about $296\,\mathrm{m}$ south of the intruder, with a covariance that claims a meter of accuracy. $K_{22}$ is negative because the intruder is almost due west of the sensor, $\varphi \approx \pi$: there $\partial y / \partial\varphi = r\cos\varphi \approx -r$, so a bearing that reads larger than predicted means a target further \emph{south}, and the gain inherits the sign.
\end{solution}

\begin{solution}{exr:ch19-banana}
Write $\varphi = \pi/2 + \epsilon$ with $\epsilon \sim \Normal(0, \sigma_\varphi^2)$, so $\cos\varphi = -\sin\epsilon$ and $\sin\varphi = \cos\epsilon$. The characteristic function of a centerd Gaussian gives $\E[\cos\epsilon] = e^{-\sigma_\varphi^2/2}$ and $\E[\sin\epsilon] = 0$; with $r$ independent of $\varphi$, $\E[r\cos\varphi] = 0$ and $\E[r\sin\varphi] = \mu_r e^{-\sigma_\varphi^2/2}$. For $\mu_r = 100$ and $\sigma_\varphi = 0.3$ this is $100\,e^{-0.045} = 95.60\,\mathrm{m}$, against the EKF's $100$. For the unscented transform, $\lambda = \alpha^2(n + \kappa) - n = 1 \cdot 3 - 2 = 1$, so $n + \lambda = 3$, $W_0^{(m)} = 1/3$ and $W_i^{(m)} = 1/6$, and $\mat{L} = \sqrt{3\mat{P}} = \operatorname{diag}(8.660, 0.5196)$. The five points $(100, \pi/2)$, $(108.660, \pi/2)$, $(100, \pi/2 \pm 0.5196)$ and $(91.340, \pi/2)$ map to $(0, 100)$, $(0, 108.660)$, $(\mp 49.66, 86.80)$ and $(0, 91.340)$. The $x$ images cancel in pairs, so the mean is $0$; the $y$ mean is $\tfrac13 \cdot 100 + \tfrac16(108.660 + 86.80 + 91.340 + 86.80) = 95.60\,\mathrm{m}$, the exact value to three decimals. Five function evaluations, no derivative, and the $4.4\,\mathrm{m}$ bias of \cref{fig:ch19-linearisation} is gone.
\end{solution}

\begin{solution}{exr:ch19-occlusion}
Hint and expected behavior. (a) Without the map likelihood nothing forbids a particle from flying through the building, so the cloud stays one connected blob straddling $x = 0$ instead of splitting; the two lobes of \cref{fig:ch19-particles} are produced by the constraint, not by the motion model. (b) Resampling at every step throws away diversity that the process noise has not yet restored: after the six blind seconds the distinct-state count collapses and the fractions west and east swing far from $68/28$, run to run, because one lobe is often wiped out entirely (sample impoverishment, \cref{sec:ch19-resampling}). (c) With $N = 200$ each lobe holds a few dozen particles, the fractions become very noisy, and $N_{\mathrm{eff}}$ at $k = 20$ often falls to single digits: the first measurement after the occlusion kills almost every particle, and the filter can lose the track altogether. Report the three numbers with at least ten different seeds; a single run tells you nothing about (b) and (c).
\end{solution}

\begin{solution}{exr:ch19-coding}
Hint. Reuse \code{simulate\_turning\_target}, \code{RangeBearingSensor} and \code{run\_tracking} from the chapter code and drive them exactly as \code{code/figures/gen\_ch19\_compare.py} does: one \code{numpy.random.Generator} per run for the truth and the measurements, a separate one for the particle filter, and the RMSE taken after the ten settling steps. (a) With $100$ runs the baseline row reproduces to about $\pm 0.05\,\mathrm{m}$; a difference of $0.02\,\mathrm{m}$ between the EKF and the UKF is inside the Monte Carlo error, which is the point of the row. (b) Expect the two curves to lie on top of each other up to about $4^\circ$ and to separate beyond it, the UKF gaining roughly $10\,\%$ by $8^\circ$; report the median as well as the mean, because a single lost run moves the mean. (c) The RMSE flattens once $N$ is a few thousand while the time per step grows linearly, so plot both on the same axis and read off the knee. (d) During a five-second occlusion the Gaussian filters coast on the coordinated-turn model and recover if the intruder did not maneuver; the particle filter spreads and keeps the alternatives, which is what \cref{sec:ch19-pf-example} shows.
\end{solution}

